\section{Analytic bounds at all embeddings}\label{sec:analysis}

Define fixed size constants
\[
 \Gamma=\max\left(1,\max_\sigma|\sigma(\gamma)|\right),\quad
 \Xi=\max\left(1,\max_\sigma|\sigma(\xi)|\right),\quad
 B_0=\Gamma(|x|+1),\quad R_0=100(1+\Gamma|x|).
\]
The maxima range over all embeddings, not just \(S\).
Thus \(B_0\ge1\), \(R_0\ge100\). Put \(R=R_0K\).

\subsection{An exact row translation}

For a polynomial column \(P\) and \(a\in\N^m\), consider the
entire function
\begin{equation}\label{eq:entire}
 f_{a,P}(z)=[u^a]P(e^z,z+u_1,\ldots,z+u_m).
\end{equation}
For \(P=Y^hX^\alpha\), this is
\(\binom{\alpha}{a}e^{hz}z^{|\alpha|-|a|}\) when \(a\le\alpha\),
and zero otherwise. In particular these functions do not depend on
the center index \(j\).

Fix \(\sigma\in S\) and put
\[
 \lambda_\sigma=\sigma(\gamma)x,\quad
 z=j\lambda_\sigma+\log(1+t),\quad
 \epsilon_i=j\sigma(\gamma)(p_i/q_i-x),\quad
 \tau_i(t)=G_i(t)-\log(1+t).
\]
The two identities
\[
 e^z=\sigma(\xi)^j(1+t),\qquad
 j\sigma(\gamma)p_i/q_i+G_i(t)+u_i
       =z+u_i+\epsilon_i+\tau_i(t)
\]
give the exact row-vector expansion
\begin{equation}\label{eq:row-expansion}
 \begin{split}
 {\rm row}_{j,\ell,\beta}
 =\sum_{\substack{a\ge\beta\\wa\le H}}\binom a\beta
   \sum_{0\le e\le a-\beta}\binom{a-\beta}{e}
       \epsilon^{a-\beta-e}
   \sum_{k=0}^{\ell}
    \left([t^k]\prod_i\tau_i(t)^{e_i}\right)
    \left([t^{\ell-k}]
       f_{a,P}(j\lambda_\sigma+\log(1+t))\right)_P .
 \end{split}
\end{equation}
The scalar multiplying a translated row is independent of the column
\(P\). This is what permits determinant multilinearity.
The identity is an equality of finite coefficient calculations:
the column space bounds \(a\), and only Taylor coefficients through
\(\ell\) occur.

\begin{lemma}[Scalar cost]\label{lem:scalar}
Every scalar in \eqref{eq:row-expansion} has absolute value at most
\begin{equation}\label{eq:scalar}
 \exp\{-\nu w(a-\beta)+H E_{\rm tr}\},
 \qquad
 E_{\rm tr}=\frac{\nu}{F_0}+\frac{\log2}{v_0}
       +\frac{\log4+\log(2K\Gamma)+\nu}{w_*}.
\end{equation}
Each row has at most
\[
 Q_H=(\lfloor H\rfloor+1)^{2m}(\lfloor H/v_0\rfloor+1)
\]
choices in this expansion.
\end{lemma}
\begin{proof}
The approximation and ceiling inequalities give
\[
 |\epsilon_i|\le K\Gamma q_i^{-\nu}
              \le K\Gamma e^\nu e^{-\nu w_i}.
\]
Each \(\tau_i\) starts at degree \(T_i\).
A nonzero coefficient in its product requires
\(\sum_i e_iT_i\le k\le\ell<H/v_0\), whence
\(we<H/F_0\).
On \(|t|=1/2\), the logarithmic tails have absolute value at most
\(\sum_{n\ge1}2^{-n}/n=\log2<1\).
Cauchy's estimate bounds the coefficient of their product by \(2^k\).
The binomial coefficients contribute at most \(4^{|a|}\), and
\(|a|\le H/w_*\).
In the approximation factor, replacing \(w(a-\beta-e)\) by
\(w(a-\beta)\) costs at most \(\nu H/F_0\).
These inequalities imply \eqref{eq:scalar}, with the harmless extra
factor \(2\) inside \(\log(2K\Gamma)\).
Finally \(w_i\ge1\), so both \(a\) and \(e\) have coordinates at
most \(\lfloor H\rfloor\); \(k\le\lfloor H/v_0\rfloor\).
\end{proof}

\subsection{The repeated-coefficient saving}

Select one translated row from each original row and omit its scalar.
Call the resulting square matrix \(T\). Let \(n_a\) count the chosen
rows with transverse index \(a\), so \(\sum_a n_a=M\).
Put
\[
 c=\frac{\log2}{4},\qquad
 E_{\rm hol}=\frac{R_0K}{w_0}+\frac{\log2}{v_0}
                   +\frac{\log(2R_0K)}{w_*}.
\]

\begin{lemma}[Collision estimate]\label{lem:collision}
Uniformly over selected columns, translated-row choices, and
\(\sigma\in S\),
\begin{equation}\label{eq:collision}
 |\det T|\le
  \exp\left\{-c\sum_a n_a^2+MH(E_{\rm hol}+\rho_H)\right\},
 \qquad \rho_H\longrightarrow0.
\end{equation}
One may take
\[
 \rho_H=
 \frac{\log M+c-\log(1-2^{-1/2})}{H}.
\]
\end{lemma}
\begin{proof}
On \(|z|\le R\), the column degree bound gives
\[
 |f_{a,P}(z)|
 \le\exp\left\{H\left(\frac R{w_0}
                        +\frac{\log(2R)}{w_*}\right)\right\}
 =:\mathcal B.
\]
On \(|t|=1/2\),
\[
 |j\lambda_\sigma+\log(1+t)|
 \le(K-1)\Gamma|x|+\log2<R/2.
\]
Expand \(f_{a,P}(z)=\sum_{d\ge0}b_{a,P,d}(z/R)^d\).
Cauchy's estimate gives \(|b_{a,P,d}|\le\mathcal B\).
The Taylor test of order \(r\) applied to \((z/R)^d\) has
absolute value at most \(2^r2^{-d}\).

If two rows use the same pair \((a,d)\), their coefficient vectors
\((b_{a,P,d})_P\) agree, and that determinant term is zero.
For a nonzero term the degrees in each \(a\)-group are distinct;
therefore
\[
 \sum_{\text{rows}}d\ge\sum_a\binom{n_a}{2}
                    =\frac{\sum_a n_a^2-M}{2}.
\]
Use half of the factor \(2^{-\sum d}\) for this lower bound and
sum its remaining half over all nonnegative degrees independently.
The latter sum is at most \((1-2^{-1/2})^{-M}\).
The determinant of each coefficient matrix is at most
\(M!\mathcal B^M\), and the sum of the Taylor-test orders is at most
\(MH/v_0\). Hence
\[
 |\det T|\le
 M!\mathcal B^M 2^{MH/v_0}
 \exp\{-c(\sum_a n_a^2-M)\}(1-2^{-1/2})^{-M}.
\]
Use \(M!\le M^M\) to obtain the asserted expression.
The series manipulations are justified by the same convergent
geometric majorant. Equation~\eqref{eq:row-count} shows
\(\log M/H\to0\). This proves the stated uniformity as well.
\end{proof}

The argument is the coefficient-collision mechanism of
\cite[Lemma 3.3]{OAIpi}. The change to real logarithmic centers
requires only the new radius bound; it does not remove the repeated
coefficient rows responsible for the saving.

\subsection{Two alternatives for the full determinant}

Fix \(0<A<1\) and \(0<\eta<1\), and let
\[
 D_A(H)=\#\{a\in\N^m:wa\le AH\},\qquad
 L_H=\frac{\eta^2M}{H D_A(H)}.
\]
The simplex counts imply
\begin{equation}\label{eq:collision-limit}
 D_A(H)\sim\frac{A^mH^m}{m!\prod_iw_i},\qquad
 L_H\longrightarrow
 L=\frac{\eta^2K\theta^m}{(m+1)v_0A^m}.
\end{equation}
If at least \(\eta M\) chosen rows have \(wa\le AH\),
Cauchy--Schwarz gives
\(\sum_a n_a^2\ge\eta^2M^2/D_A(H)=MH L_H\).
For a nonzero scalar, \(a\ge\beta\); thus the approximation
exponent in \eqref{eq:scalar} is nonpositive and may be discarded.
If fewer than \(\eta M\) rows have \(wa\le AH\), then
\[
 \sum_{\rm rows}w(a-\beta)>
 MH\bigl(A(1-\eta)-\bar b_H\bigr).
\]
In this alternative discard the nonpositive collision exponent.
The determinant expansion has at most \(Q_H^M\) terms.
Taking logarithms of the triangle inequality yields, at every
compatible embedding,
\begin{equation}\label{eq:good-bound}
 \frac{\log|\sigma(\Delta_H)|}{MH}
 \le E_{\rm tr}+E_{\rm hol}
     +\rho_H+\frac{\log Q_H}{H}
     +\max\{-cL_H,-\nu(A(1-\eta)-\bar b_H)\}.
\end{equation}
This upper bound applies to every selected nonzero full-row minor.
It is uniform in its selected columns; no analytic choice of a special
minor is made after the arithmetic choice.

\subsection{The other embeddings}

For an arbitrary embedding \(\sigma\), use
\eqref{eq:binomial-entry} directly. On \(|t|=1/2\),
\[
 |j\sigma(\gamma)p_i/q_i|\le KB_0,\qquad |G_i(t)|<1.
\]
The binomial coefficient costs at most \(2^{|\alpha|}\);
\(|\sigma(\xi)^{jh}|\le\Xi^{Kh}\);
\(|1+t|^h\le(3/2)^h\); and coefficient extraction costs \(2^\ell\).
It follows that every entry is at most \(\exp(HU)\), where the
convenient common bound used by the formal proof is
\begin{equation}\label{eq:direct-error}
 U=\frac{(\log\Xi+\log(3/2))K}{w_0}
       +\frac{\log(2(KB_0+2))}{w_*}
       +\frac{\log2}{v_0}\ge0.
\end{equation}
Thus
\begin{equation}\label{eq:other-bound}
 \frac{\log|\sigma(\Delta_H)|}{MH}\le U+\frac{\log M}{H}.
\end{equation}
We claim no collision saving at incompatible embeddings.

\begin{proposition}[All-embedding upper bound]\label{prop:analytic}
Let
\[
 E_{\rm an}=E_{\rm tr}+E_{\rm hol}+\kappa U.
\]
For the same minor used in Proposition~\ref{prop:arithmetic},
\begin{equation}\label{eq:analytic}
 \mathcal N_H\le E_{\rm an}+o(1)
       +\max\{-cL_H,-\nu(A(1-\eta)-\bar b_H)\}.
\end{equation}
The remainder tends to zero for the fixed data and is uniform over
the selected columns.
\end{proposition}
\begin{proof}
Sum \eqref{eq:good-bound} over \(S\), and
\eqref{eq:other-bound} over its complement, then divide by \(s\).
Since \(U+\log M/H\ge0\), replacing the complement size by \(d\)
only enlarges the upper bound.
One resulting remainder is
\[
 \rho_H+\frac{\log Q_H}{H}+\kappa\frac{\log M}{H},
\]
which tends to zero by the polynomial growth of \(M\) and \(Q_H\).
\end{proof}
